% ============================================================
\chapter{Measures}
\label{ch:measures}
% ============================================================

Now that we have $\sigma$-algebras --- the collections of sets to which we
wish to assign a size --- it is time to define that ``size'' itself. This is
the role of the notion of \emph{measure}, which simultaneously generalizes
length, area, volume, and probability. This unification, due to Lebesgue and
refined by Carath\'eodory, is one of the great achievements of modern
mathematics.

% ------------------------------------------------------------
\section{Definition and first properties}
% ------------------------------------------------------------

\begin{definition}[Measure]
Let $(X, \mathcal{A})$ be a measurable space. A \emph{measure} on
$(X, \mathcal{A})$ is a function $\mu : \mathcal{A} \to [0, +\infty]$ satisfying:
\begin{enumerate}
    \item[(i)] $\mu(\varnothing) = 0$;
    \item[(ii)] ($\sigma$-additivity) for every sequence $(A_n)_{n \in \N}$ of
    \emph{pairwise disjoint} elements of $\mathcal{A}$:
    \[
    \mu\!\left(\bigsqcup_{n=0}^{\infty} A_n\right) = \sum_{n=0}^{\infty} \mu(A_n).
    \]
\end{enumerate}
The triple $(X, \mathcal{A}, \mu)$ is called a \emph{measure space}.
\end{definition}

\begin{definition}[Types of measures]
Let $(X, \mathcal{A}, \mu)$ be a measure space.
\begin{itemize}
    \item $\mu$ is \emph{finite} if $\mu(X) < \infty$.
    \item $\mu$ is a \emph{probability measure} if $\mu(X) = 1$.
    \item $\mu$ is \emph{$\sigma$-finite} if there exists a sequence
    $(X_n)_{n \in \N} \subset \mathcal{A}$ with $X = \bigcup_n X_n$ and
    $\mu(X_n) < \infty$ for all $n$.
    \item $\mu$ is \emph{complete} if for every $N \in \mathcal{A}$ with
    $\mu(N) = 0$ and every $A \subset N$, we have $A \in \mathcal{A}$.
\end{itemize}
\end{definition}

\begin{proposition}[Monotonicity and subadditivity]
\label{prop:monotonicity}
Let $(X, \mathcal{A}, \mu)$ be a measure space.
\begin{enumerate}
    \item \textbf{Monotonicity}: if $A \subset B$ with $A, B \in \mathcal{A}$, then
    $\mu(A) \leq \mu(B)$.
    \item \textbf{Excision}: if $A \subset B$ and $\mu(A) < \infty$, then
    $\mu(B \setminus A) = \mu(B) - \mu(A)$.
    \item \textbf{Countable subadditivity}: for any sequence
    $(A_n) \subset \mathcal{A}$,
    $\mu\!\left(\bigcup_n A_n\right) \leq \sum_n \mu(A_n)$.
\end{enumerate}
\end{proposition}

\begin{proof}
(1) Write $B = A \sqcup (B \setminus A)$, so
$\mu(B) = \mu(A) + \mu(B \setminus A) \geq \mu(A)$.

(2) Follows directly from the decomposition $B = A \sqcup (B \setminus A)$ and
the finiteness of $\mu(A)$.

(3) Set $B_0 = A_0$ and $B_n = A_n \setminus \bigcup_{k=0}^{n-1} A_k$ for
$n \geq 1$. Then $(B_n)$ is a disjoint sequence with $\bigcup_n B_n =
\bigcup_n A_n$ and $B_n \subset A_n$. Therefore
\[
\mu\!\left(\bigcup_n A_n\right) = \mu\!\left(\bigsqcup_n B_n\right)
= \sum_n \mu(B_n) \leq \sum_n \mu(A_n). \qedhere
\]
\end{proof}

\section{Continuity of measures}

\begin{theorem}[Monotone continuity]
\label{thm:monotone_continuity}
Let $(X, \mathcal{A}, \mu)$ be a measure space.
\begin{enumerate}
    \item \textbf{Continuity from below}: if $A_n \uparrow A$ (increasing sequence),
    then $\mu(A_n) \uparrow \mu(A)$.
    \item \textbf{Continuity from above}: if $A_n \downarrow A$ (decreasing sequence)
    and $\mu(A_1) < \infty$, then $\mu(A_n) \downarrow \mu(A)$.
\end{enumerate}
\end{theorem}

\begin{proof}
(1) Set $B_0 = A_0$ and $B_n = A_n \setminus A_{n-1}$ for $n \geq 1$. Then
$(B_n)$ is disjoint with $\bigsqcup_{k=0}^n B_k = A_n$ and
$\bigsqcup_{k=0}^\infty B_k = A$. By $\sigma$-additivity:
\[
\mu(A) = \sum_{k=0}^{\infty} \mu(B_k) = \lim_{n \to \infty} \sum_{k=0}^n \mu(B_k)
= \lim_{n \to \infty} \mu(A_n).
\]

(2) Set $C_n = A_1 \setminus A_n$. Then $C_n \uparrow A_1 \setminus A$. By (1),
$\mu(C_n) \uparrow \mu(A_1 \setminus A)$. Since $\mu(A_1) < \infty$:
$\mu(C_n) = \mu(A_1) - \mu(A_n)$ and $\mu(A_1 \setminus A) = \mu(A_1) - \mu(A)$,
giving $\mu(A_n) \downarrow \mu(A)$.
\end{proof}

\begin{warning}[title=\textbf{Finiteness condition}]
The hypothesis $\mu(A_1) < \infty$ in (2) is essential. Counterexample:
$A_n = [n, +\infty)$ with Lebesgue measure. Then $A_n \downarrow \varnothing$ but
$\mu(A_n) = +\infty$ for all $n$.
\end{warning}

\begin{lemma}[Borel--Cantelli lemma]
\label{lem:borel_cantelli}
Let $(X, \mathcal{A}, \mu)$ be a measure space and $(A_n) \subset \mathcal{A}$. If
$\sum_{n=0}^{\infty} \mu(A_n) < \infty$, then
\[
\mu\!\left(\limsup_{n \to \infty} A_n\right) = 0,
\]
where $\limsup_n A_n = \bigcap_{n=0}^{\infty} \bigcup_{k=n}^{\infty} A_k$ is the
set of points belonging to infinitely many $A_n$.
\end{lemma}

\begin{proof}
Set $B_n = \bigcup_{k=n}^{\infty} A_k$. Then $B_n \downarrow \limsup_n A_n$ and
\[
\mu(B_n) \leq \sum_{k=n}^{\infty} \mu(A_k) \to 0 \quad (n \to \infty)
\]
since the tail of a convergent series tends to zero. Since $\mu(B_0) \leq
\sum_k \mu(A_k) < \infty$, continuity from above gives
$\mu(\limsup_n A_n) = \lim_n \mu(B_n) = 0$.
\end{proof}

\section{Fundamental examples}

\begin{example}[Counting measure]
Let $X$ be any set and $\mathcal{A} = \mathcal{P}(X)$. The \emph{counting measure}
is defined by:
\[
\mu(A) = \begin{cases} \abs{A} & \text{if } A \text{ is finite,} \\
+\infty & \text{otherwise.} \end{cases}
\]
It is a measure on $(X, \mathcal{P}(X))$, $\sigma$-finite if and only if $X$ is
countable.
\end{example}

\begin{example}[Dirac measure]
Let $x_0 \in X$. The \emph{Dirac measure} (or point mass) at $x_0$ is the
probability measure:
\[
\delta_{x_0}(A) = \begin{cases} 1 & \text{if } x_0 \in A, \\ 0 & \text{if }
x_0 \notin A. \end{cases}
\]
\end{example}

\begin{example}[Linear combinations]
If $\mu_1, \mu_2$ are measures on $(X, \mathcal{A})$ and $\alpha, \beta \geq 0$,
then $\alpha \mu_1 + \beta \mu_2$ is a measure. More generally, if
$(x_n)_{n \in \N} \subset X$ and $(a_n)_{n \in \N} \subset [0, +\infty)$, then
$\mu = \sum_n a_n \delta_{x_n}$ is a measure.
\end{example}

\section{Pre-measures and extension}

\begin{definition}[Pre-measure]
Let $\mathcal{A}_0$ be an algebra on $X$. A function
$\mu_0 : \mathcal{A}_0 \to [0, +\infty]$ is a \emph{pre-measure} if:
\begin{enumerate}
    \item[(i)] $\mu_0(\varnothing) = 0$;
    \item[(ii)] for every disjoint sequence $(A_n) \subset \mathcal{A}_0$ such that
    $\bigsqcup_n A_n \in \mathcal{A}_0$:
    $\mu_0\!\left(\bigsqcup_n A_n\right) = \sum_n \mu_0(A_n)$.
\end{enumerate}
\end{definition}

\begin{remark}
The difference from a measure is that $\mathcal{A}_0$ is only an algebra, and
$\sigma$-additivity is required only when the union falls in $\mathcal{A}_0$.
\end{remark}

\begin{proposition}[Criterion for $\sigma$-additivity]
\label{prop:sigma_add_criterion}
Let $\mu_0$ be a finitely additive set function on an algebra $\mathcal{A}_0$ with
$\mu_0(\varnothing) = 0$. The following are equivalent:
\begin{enumerate}
    \item $\mu_0$ is a pre-measure ($\sigma$-additive).
    \item $\mu_0$ is continuous from below: if $(A_n) \subset \mathcal{A}_0$,
    $A_n \uparrow A \in \mathcal{A}_0$, then $\mu_0(A_n) \uparrow \mu_0(A)$.
    \item $\mu_0$ is continuous at $\varnothing$: if $(A_n) \subset \mathcal{A}_0$,
    $A_n \downarrow \varnothing$, then $\mu_0(A_n) \downarrow 0$.
\end{enumerate}
\end{proposition}

\begin{proof}
$(1) \Rightarrow (2)$: same proof as Theorem~\ref{thm:monotone_continuity}(1).

$(2) \Rightarrow (3)$: if $A_n \downarrow \varnothing$ and
$\mu_0(A_1) < \infty$, apply the same technique as
Theorem~\ref{thm:monotone_continuity}(2). If $\mu_0(A_1) = +\infty$, the result is
trivial.

$(3) \Rightarrow (1)$: let $(A_n)$ be disjoint with $A = \bigsqcup_n A_n \in
\mathcal{A}_0$. Set $B_N = A \setminus \bigsqcup_{n=0}^{N} A_n$. By finite
additivity, $\mu_0(A) = \sum_{n=0}^{N} \mu_0(A_n) + \mu_0(B_N)$. Since
$B_N \downarrow \varnothing$, we get $\mu_0(B_N) \to 0$, hence
$\mu_0(A) = \sum_n \mu_0(A_n)$.
\end{proof}

\section{Extension theorem}

\begin{theorem}[Carath\'eodory extension --- preliminary statement]
\label{thm:extension_preliminary}
Let $\mu_0$ be a pre-measure on an algebra $\mathcal{A}_0$. Then $\mu_0$ extends
to a measure $\mu$ on $\sigma(\mathcal{A}_0)$. If $\mu_0$ is $\sigma$-finite, this
extension is unique.
\end{theorem}

The complete proof will be given in Chapter~\ref{ch:caratheodory}, but we sketch
the construction here.

\begin{keyformulas}[title=\textbf{Outer measure construction}]
One defines the \emph{outer measure} associated with $\mu_0$ by:
\[
\mu^*(A) = \inf\left\{\sum_{n=0}^{\infty} \mu_0(A_n) :
(A_n) \subset \mathcal{A}_0,\, A \subset \bigcup_{n=0}^{\infty} A_n\right\}
\]
for every $A \subset X$. Then one restricts $\mu^*$ to the
\emph{Carath\'eodory-measurable} sets.
\end{keyformulas}

\section{Completion of a measure}

\begin{definition}[Completion]
Let $(X, \mathcal{A}, \mu)$ be a measure space. The \emph{completion} of $\mu$ is
the measure $\overline{\mu}$ defined on the $\sigma$-algebra:
\[
\overline{\mathcal{A}} = \{A \cup N : A \in \mathcal{A},\, N \subset M \text{ for
some } M \in \mathcal{A} \text{ with } \mu(M) = 0\}
\]
by $\overline{\mu}(A \cup N) = \mu(A)$.
\end{definition}

\begin{proposition}
$\overline{\mathcal{A}}$ is a $\sigma$-algebra and $\overline{\mu}$ is a complete
measure extending $\mu$.
\end{proposition}

\begin{proof}
We verify that $\overline{\mathcal{A}}$ is a $\sigma$-algebra.

\emph{Closure under complementation}: if $E = A \cup N$ with $N \subset M$,
$\mu(M) = 0$, then $E^c = A^c \cap N^c = (A^c \cap M^c) \cup (A^c \cap M
\setminus N)$. Now $A^c \cap M^c \in \mathcal{A}$ and $A^c \cap M \setminus N
\subset M$, so $E^c \in \overline{\mathcal{A}}$.

\emph{Closure under countable unions}: if $E_n = A_n \cup N_n$ with
$N_n \subset M_n$, $\mu(M_n) = 0$, then $\bigcup_n E_n = \bigcup_n A_n \cup
\bigcup_n N_n$ with $\bigcup_n N_n \subset \bigcup_n M_n$ and
$\mu(\bigcup_n M_n) \leq \sum_n \mu(M_n) = 0$.

Well-definedness of $\overline{\mu}$ (independence of the decomposition) is verified
directly, and $\sigma$-additivity follows from that of $\mu$.
\end{proof}

\section{Image measure}

\begin{definition}[Image measure]
Let $(X, \mathcal{A}, \mu)$ be a measure space and $f : X \to Y$ a measurable map
to $(Y, \mathcal{B})$. The \emph{image measure} (or \emph{push-forward}) of $\mu$
under $f$ is the measure $f_*\mu$ on $(Y, \mathcal{B})$ defined by:
\[
(f_*\mu)(B) = \mu(f^{-1}(B)), \quad B \in \mathcal{B}.
\]
\end{definition}

\begin{proposition}
$f_*\mu$ is indeed a measure on $(Y, \mathcal{B})$.
\end{proposition}

\begin{proof}
$(f_*\mu)(\varnothing) = \mu(f^{-1}(\varnothing)) = \mu(\varnothing) = 0$. For a
disjoint sequence $(B_n)$ in $\mathcal{B}$, $(f^{-1}(B_n))$ is disjoint in
$\mathcal{A}$ and $f^{-1}(\bigsqcup_n B_n) = \bigsqcup_n f^{-1}(B_n)$, so
$(f_*\mu)(\bigsqcup_n B_n) = \mu(\bigsqcup_n f^{-1}(B_n)) = \sum_n
\mu(f^{-1}(B_n)) = \sum_n (f_*\mu)(B_n)$.
\end{proof}

\section{Exercises}

\begin{exercise}
Let $(X, \mathcal{A}, \mu)$ be a measure space and $(A_n) \subset \mathcal{A}$.
Show that
$\mu(\liminf_n A_n) \leq \liminf_n \mu(A_n)$
(Fatou's lemma for sets).
\end{exercise}

\begin{exercise}
Let $\mu$ be a measure on $(\R, \mathcal{B}(\R))$ such that $\mu(K) < \infty$ for
every compact $K$. Show that $\mu$ is $\sigma$-finite.
\end{exercise}

\begin{exercise}
Show that the counting measure on $(\R, \mathcal{P}(\R))$ is not $\sigma$-finite.
\end{exercise}

\begin{exercise}
Let $(X, \mathcal{A}, \mu)$ be a measure space with $\mu(X) < \infty$ and
$(A_n) \subset \mathcal{A}$ a sequence with $\mu(A_n) \geq \alpha > 0$ for all $n$.
Show that there exists $x \in X$ belonging to infinitely many $A_n$.
\end{exercise}

\begin{exercise}
Let $\mu$ be a finite measure on $(\N, \mathcal{P}(\N))$. Show that $\mu$ is
determined by the values $\mu(\{n\})$, $n \in \N$, and that
$\mu = \sum_{n=0}^{\infty} \mu(\{n\}) \delta_n$.
\end{exercise}

\begin{exercise}[Lebesgue--Stieltjes measure]
Let $F : \R \to \R$ be an increasing right-continuous function. Define
$\mu_0((a, b]) = F(b) - F(a)$ on the algebra $\mathcal{A}_0$ of finite disjoint
unions of half-open intervals $(a, b]$.
\begin{enumerate}
    \item Show that $\mu_0$ is a pre-measure on $\mathcal{A}_0$.
    \item Deduce the existence of a unique measure $\mu_F$ on $\mathcal{B}(\R)$
    such that $\mu_F((a, b]) = F(b) - F(a)$.
\end{enumerate}
\end{exercise}
